\chapter{LITERATURE REVIEW}

\hspace{\parindent} 
Technology has become an essential tool in addressing the challenges associated with an aging population. One area of growing interest is the development of smart companion systems aimed at supporting elderly individuals in their daily lives. These systems integrate artificial intelligence and communication platforms to enhance quality of life, improve safety, and foster independence. This research aims to provide personalized and accessible care for elderly individuals and their families, supporting their emotional well-being.
This literature review investigates the application of each selected study to the development of a smart companion system for elderly care. The primary goals of the study are to implement a function to predict emotions based on audio and facial recognition data, conduct natural conversations, introduce a video call function for real-time communication, and develop a monitoring system that functions as a camera for elderly care. These features aim to meet both the emotional and practical needs of elderly individuals, offering comprehensive support.
The scope of this research is shaped by the areas explored in the selected studies. The Multimodal Emotion Recognition Using Feature Fusion: An LLM-Based Approach focuses on emotion recognition using multimodal inputs by converting audio and visual features into structured text, offering a new approach for LLM-based classification accuracy. The AI Affective Conversational Robot with Hybrid Generative-Based and Retrieval-Based Dialogue Models investigates emotionally-aware dialogue systems through deep learning-based sentiment analysis and dialogue generation. The Development of an Intelligent Dialogue Agent for Older Adults explores how reinforcement learning techniques can adapt conversation topics based on user preferences to promote social companionship. These studies collectively provide insights into multimodal emotion recognition, sentiment prediction, and personalized dialogue generation, forming the foundation for this research.


\setcounter{chapter}{2}
\section[Multimodal Emotion Recognition Using Feature Fusion: An LLM-Based Approach]{Multimodal Emotion Recognition Using Feature Fusion: An LLM-Based Approach \cite{paper_1}}

\hspace*{17pt}Multimodal emotion recognition has evolved from traditional machine learning methods to deep learning-based multimodal fusion techniques, integrating speech, facial expressions, and textual data. While deep learning models like CNNs, RNNs, and transformers have improved classification accuracy, they often suffer from high computational costs and feature alignment challenges. Recent advancements in Large Language Models (LLMs) such as BERT, RoBERTa, and GPT-3 have enabled emotion recognition through textual representations of multimodal inputs. This paper addresses key gaps by proposing a novel feature fusion approach that converts audio and visual features into structured text, processed by a fine-tuned DistilRoBERTa model for emotion classification. The method achieves state-of-the-art accuracy on benchmark datasets, demonstrating that LLM-based textual fusion can match or exceed traditional multimodal deep learning models while improving scalability and interpretability. This approach bridges the gap between multimodal emotion recognition and NLP, offering a promising direction for future affective computing applications.
\


\newpage
\section[AI Affective Conversational Robot with Hybrid Generative-Based and Retrieval-Based Dialogue Models]{AI Affective Conversational Robot with Hybrid Generative-Based and Retrieval-Based Dialogue Models \cite{paper_3}}
\hspace*{17pt} ChatBot technology has become a widely used in various application fields. An important topic in the research on conversational robots is the improvement of their temperature during operation for enhanced user interaction. In this study, we propose an artificial intelligence affective conversational robot (AIACR), which is an integration of an artificial intelligence deep learning sentiment analysis model and generative-and retrieval-based dialogue models. The sentiment analysis model developed in this study uses three models, namely, multilayer perceptron (MLP), long short-term memory (LSTM) and bidirectional long short-term memory (BiLSTM). Moreover, word2vec and semantics are utilized as the basis for similarity ranking models. The deep learning dialogue model, sentiment analysis model, and similarity model were integrated and compared as well. The experimental results show that the sentiment analysis model, similarity model, and dialogue model respectively utilize BiLSTM, word2vec, and the retrieval-based model to achieve the best dialogue performance. The major research contributions of this study are the developed AIACR and the proposed affective conversational robot index (ACR Index) as a criterion for evaluating the effectiveness of emotional dialogue robots.\

The main objectives of this paper are as follows:
\begin{itemize}
	\item To compare the effects of conversational robots in retrieval- and generative-based models. 
	\item To use artificial intelligence (AI) technology such as multilayer perceptron (MLP), long short-term memory (LSTM), and bidirectional LSTM (BiLSTM) to conduct sentiment analysis comparison. 
	\item To integrate the dialogue model with the sentiment prediction model to produce a dialogue with emotions.
\end{itemize}



	\newpage
	\section[Development of an Intelligent Dialogue Agent for Older Adults : Evaluation of Functions to Control Spontaneous Talk and Coordinate Speech Content]{ Development of an Intelligent Dialogue Agent for Older Adults : Evaluation of Functions to Control Spontaneous Talk and Coordinate Speech Content \cite{paper_2}}
	\hspace*{17pt} A Reinforcement Learning-Based Approach Intelligent Dialogue Agents (IDAs) have been increasingly utilized to support elderly care by offering social companionship and promoting preventive care activities. Traditional voice-interaction devices like Google Home and Amazon Echo provide information in response to user queries but cannot initiate spontaneous conversations. Recent research has focused on developing IDAs with spontaneous talking functions (STF) to foster more natural and friendly interactions, enhancing users' sense of familiarity and engagement.
	Kitakoshi et al. (2019) proposed an IDA that combines STF with a Speech-Content Coordinating Function (SCCF) to adapt speech content based on user preferences and behavioral trends. The SCCF employs reinforcement learning (RL) algorithms, specifically a policy gradient method, to select appropriate speech topics at the right moments. Experiments conducted with students demonstrated that users responded positively to personalized content but highlighted the need to optimize speech length and RL parameters to improve learning performance. The study suggests that integrating adaptive dialogue generation with preventive care systems can encourage habitual use among older adults, contributing to comprehensive preventive care frameworks.	\

\newpage
\section[A Survey on LLM-as-a-Judge]{A Survey on LLM-as-a-Judge \cite{paper_6}}
\hspace*{17pt} The rapid proliferation of Large Language Models (LLMs) has necessitated scalable and reliable evaluation frameworks, especially for open-ended generation tasks where traditional metrics like BLEU or ROUGE fall short. This paper presents a comprehensive survey of the "LLM-as-a-Judge" paradigm, which leverages advanced LLMs (e.g., GPT-4, Gemini) to evaluate the quality of outputs generated by other models. The authors categorize existing approaches into three primary methodologies: \textbf{Pairwise Comparison}, where the judge selects the better of two responses; \textbf{Score-Based Evaluation}, where the judge assigns a numerical rating based on specific criteria (e.g., helpfulness, safety); and \textbf{Reference-Based Evaluation}, which compares the output against a gold standard.

The survey highlights that while LLM-based evaluation offers significant advantages in terms of speed, cost, and reproducibility compared to human evaluation, it is not without challenges. Key biases such as \textbf{position bias} (preferring the first option), \textbf{verbosity bias} (favoring longer responses), and \textbf{self-preference bias} (preferring outputs from the same model family) are critically analyzed. The paper concludes that with careful prompt engineering and bias mitigation strategies, LLM-as-a-Judge can achieve high correlation with human judgment, making it a viable standard for modern AI development.

This research is directly applicable to the development of the "Buddy" system, validating the use of an automated evaluation pipeline to assess the dialogue quality of the edge-deployed companion model. By adopting the score-based evaluation methodology described, this project ensures a rigorous and objective assessment of the chatbot's empathetic capabilities without the prohibitive resource costs of extensive human trials.

\newpage
\section[AI-Enhanced Social Robots for Older Adults Care: Evaluating the Efficacy of ChatGPT-Powered Storytelling in the EBO Platform]{AI-Enhanced Social Robots for Older Adults Care: Evaluating the Efficacy of ChatGPT-Powered Storytelling in the EBO Platform \cite{paper_7}}
\hspace*{17pt} Social robots are increasingly recognized as potent therapeutic tools for addressing the cognitive and emotional needs of the aging population. However, traditional social robots often rely on pre-scripted interactions, which can limit long-term engagement and adaptability. This study addresses these limitations by introducing an enhanced version of the EBO social robot, seamlessly integrated with the CORTEX cognitive architecture and powered by a Large Language Model (LLM), specifically ChatGPT. The core innovation lies in the system's ability to generate dynamic, natural-sounding narratives that are not only engaging but also contextually tailored to the user's cognitive profile and personal interests.

A key feature of the proposed framework is the "Therapist-in-the-Loop" design, which allows clinical professionals to guide the AI's storytelling parameters, ensuring therapeutic relevance while leveraging the LLM's generative capabilities for variety. The study validates this approach through a case study involving older adults and therapists. Results indicate that the AI-enhanced EBO robot significantly improves user engagement compared to conventional methods, effectively facilitating interactive storytelling sessions that promote cognitive stimulation and socio-emotional connection. This research provides a crucial reference for the "Buddy" project, demonstrating that integrating generative AI with social robotics can transform passive entertainment into active, personalized care, thereby validating the architectural decision to employ GEMMA for dynamic dialogue generation in our own system.

\newpage
\section[Thonburian Whisper: Robust Fine-tuned and Distilled Whisper for Thai]{Thonburian Whisper: Robust Fine-tuned and Distilled Whisper for Thai \cite{aung-etal-2024-thonburian}}
\hspace*{17pt} Automatic Speech Recognition (ASR) for tonal languages like Thai presents significant challenges due to the scarcity of large-scale, high-quality labeled datasets compared to English. While OpenAI's Whisper model has demonstrated impressive zero-shot capabilities, its performance on specific low-resource languages can be inconsistent, particularly in noisy environments. This study addresses these limitations by introducing \textbf{Thonburian Whisper}, a specialized adaptation of the Whisper architecture that undergoes rigorous fine-tuning and knowledge distillation specifically for the Thai language. The authors leverage a composite dataset comprising Common Voice 11.0, Google Fleurs, and proprietary curation to enhance the model's robustness against diverse acoustic conditions.

A critical contribution of this work is the application of \textbf{model distillation}, where a larger "teacher" model (Whisper `large-v2`) transfers knowledge to a smaller, more efficient "student" model. This process significantly reduces the computational footprint required for inference without a proportional loss in accuracy, making it feasible to deploy transformer-based ASR on resource-constrained edge devices. The study establishes a new state-of-the-art benchmark for Thai ASR, achieving a Word Error Rate (WER) of \textbf{11.01\%} on the Common Voice test set. For the "Buddy" project, this research serves as a foundational reference, justifying the selection of the Faster Whisper architecture and providing a validated baseline metric against which our system's STT performance is evaluated.

\newpage
\section[Emotional design: Creating user interfaces that evoke emotions]{Emotional design: Creating user interfaces that evoke emotions \cite{paper_8}}
\hspace*{17pt} While advanced AI algorithms drive the functional intelligence of companion systems, the User Interface (UI) design plays a pivotal role in determining user acceptance and emotional connection. This thesis explores the concept of "Emotional Design" in mobile applications, investigating how visual elements—such as color palettes, typography, and layout—can be strategically engineered to evoke specific positive emotions and enhance user satisfaction. The study is grounded in Donald Norman’s three levels of emotional design: \textbf{Visceral} (immediate aesthetic appeal), \textbf{Behavioral} (usability and functionality), and \textbf{Reflective} (long-term meaning and personal connection).

Through comprehensive user surveys and case studies across social and health industries, the research demonstrates that emotional design elements significantly influence user engagement and loyalty. For the "Buddy" project, specifically the Caregiver and Elderly applications, these findings underscore the importance of a UI that is not just functional but also emotionally resonant. Applying these principles—for instance, using calming colors for the elderly interface to reduce anxiety or clear, reassuring typography for health alerts—can bridge the gap between complex AI technology and the sensitive needs of elderly users, fostering a deeper sense of trust and comfort.

\newpage
\section[Performance Evaluation and Application Potential of Small Large Language Models in Complex Sentiment Analysis Tasks]{Performance Evaluation and Application Potential of Small Large Language Models in Complex Sentiment Analysis Tasks \cite{paper_9}}
\hspace*{17pt} The deployment of Large Language Models (LLMs) for sensitive applications, such as healthcare and elderly care, often faces challenges related to data privacy, latency, and operational costs. While giant models like GPT-4 set the benchmark for performance, their closed-source nature and cloud dependency make them less ideal for privacy-critical edge devices. This paper investigates the viability of "Small Large Language Models" (sLLMs)—models with significantly fewer parameters (e.g., 1.5B to 7B)—as robust alternatives for complex sentiment analysis tasks. The authors conducted a comparative evaluation of several sLLMs against GPT-3.5 and GPT-4, focusing on accuracy, instruction-following capability, and resource consumption (VRAM and generation time).

The findings demonstrate that modern sLLMs, such as Qwen2.5 (similar in scale to the GEMMA models used in this project), can achieve sentiment analysis F1-scores comparable to or even exceeding GPT-3.5 in specialized domains. The study highlights that with appropriate prompting strategies, these smaller models offer a superior trade-off between performance and efficiency, making them highly suitable for resource-constrained environments. This research directly supports the architectural choices of the "Buddy" system, validating the use of the lightweight \textbf{GEMMA 3:1b} and \textbf{4b} models for local, privacy-preserving emotion recognition and dialogue generation, ensuring that elderly user data remains secure on the device while maintaining high analytical accuracy.

\newpage
\section[Personality Perceptions of Conversational Agents: A Task-Based Analysis Using Thai as the Conversational Language]{Personality Perceptions of Conversational Agents: A Task-Based Analysis Using Thai as the Conversational Language \cite{paper_10}}
\hspace*{17pt} The perceived personality of a Conversational Agent (CA) significantly influences user trust and engagement, particularly in cultures where social hierarchy and linguistic politeness are deeply ingrained. This study specifically investigates how Thai users perceive CA personalities and how these perceptions affect task performance across social and functional domains. The authors propose a 7-dimensional CA personality model, identifying key traits such as "Maturity-Juvenility" and "Intelligence-Ineptness" as critical drivers of user acceptance.

The findings reveal that for social tasks—such as the companionship provided by the "Buddy" system—users are most influenced by the "Maturity-Juvenility" dimension. A CA that exhibits a consistent, culturally appropriate persona (e.g., respectful but approachable) fosters deeper social connection. This research is instrumental for the "Buddy" project, validating the decision to fine-tune the system prompt ($S$) to exclude overly formal particles (`Ka/Krub`) in favor of a "family-like" tone. It underscores that for Thai elderly users, the *style* of communication is as important as the content, and aligning the bot's personality with the user's expectations of a "grandchild" or "caregiver" role is essential for long-term adoption.

\newpage
\section[Conversational Affective Social Robots for Ageing and Dementia Support]{Conversational Affective Social Robots for Ageing and Dementia Support \cite{paper_11}}
\hspace*{17pt} Socially Assistive Robots (SARs) represent a promising frontier for supporting aging populations, particularly those with dementia, by enhancing mental health and enabling independence at home. However, despite significant research growth, the field faces persistent challenges in establishing long-term trust and achieving meaningful "clinical translation"—the move from lab prototypes to real-world health benefits. This paper reviews two decades of SAR research, identifying a critical gap: the "immaturity" of affective human-robot interactions. Existing systems often fail to sustain engagement because they lack robust conversational abilities and emotional intelligence.

The authors argue that the deployment of robots with advanced conversational capabilities is fundamental for robustness. They propose a roadmap for future development that emphasizes user-centered design, specifically the ability to handle communication breakdowns—a frequent occurrence with elderly users who may have cognitive impairments. For the "Buddy" project, this review serves as both a validation of the problem space and a guide for solution design. By integrating a sophisticated Large Language Model (GEMMA) capable of context-aware dialogue and "repairing" misunderstandings, alongside an emotion recognition module, "Buddy" directly addresses the "affective gap" identified in the literature, positioning itself as a next-generation tool for longitudinal health monitoring and companionship.
